\documentclass[11pt]{article}
\usepackage{preamble}
\setlength{\parskip}{4pt}

\begin{document}

\daybanner{AP Statistics \textperiodcentered\ Unit 3 \textperiodcentered\ Topic 3.14 \textperiodcentered\ B: 2026-12-22 \textperiodcentered\ E: 2027-02-05 \textperiodcentered\ TEACHER KEY}{One Table, Two Study Designs}

\sectionbanner{CED TETHER}

{\small
\begin{itemize}[leftmargin=1.5em,itemsep=2pt,topsep=2pt]
  \item \textbf{Skill 1.F} --- Identify null and alternative hypotheses.
  \item \textbf{Skill 1.E} --- Identify an appropriate inference method for significance tests.
  \item \textbf{Skill 4.C} --- Verify that inference procedures apply in a given situation.
  \item \textbf{EU VAR-8} --- The chi-square distribution may be used to model variation.
  \item \textbf{LO VAR-8.I} --- Identify the null and alternative hypotheses for a chi-square test for homogeneity or independence.
  \item \textbf{EK VAR-8.I.1} --- The appropriate hypotheses for a chi-square test for homogeneity are: $H_0$: There is no difference in distributions of a categorical variable across populations or treatments. $H_a$: There is a difference in distributions of a categorical variable across populations or treatments.
  \item \textbf{EK VAR-8.I.2} --- The appropriate hypotheses for a chi-square test for independence are: $H_0$: There is no association between two categorical variables in a given population (or the two variables are independent). $H_a$: Two categorical variables in a population are associated or dependent.
  \item \textbf{LO VAR-8.J} --- Identify an appropriate testing method for comparing distributions in two-way tables of categorical data.
  \item \textbf{EK VAR-8.J.1} --- When comparing distributions to determine whether proportions in each category for categorical data collected from different populations are the same, the appropriate test is the chi-square test for homogeneity.
  \item \textbf{EK VAR-8.J.2} --- To determine whether row and column variables in a two-way table of categorical data might be associated in the population from which the data were sampled, the appropriate test is the chi-square test for independence.
  \item \textbf{LO VAR-8.K} --- Verify the conditions for making statistical inferences when testing a chi-square distribution for independence or homogeneity.
  \item \textbf{EK VAR-8.K.1} --- Conditions for chi-square tests on two-way tables (homogeneity or independence): a. Independence check: i. For a test for independence: data should be collected using a simple random sample.
  \item ii. For a test for homogeneity: data should be collected using a stratified random sample or randomized experiment.
  \item iii. When sampling without replacement, check that n $\leq$ 10\% N.
  \item b. Large counts (shape): all expected counts should be greater than 5.
\end{itemize}
}
\textbf{Essential question:} Why can identical counts require different chi-square procedure names and hypotheses?\par
\textbf{DOK-3 skill:} 4.C \textperiodcentered\ \textbf{FRQ pattern:} {\small\texttt{design-\allowbreak{}controls-\allowbreak{}chi-\allowbreak{}square-\allowbreak{}procedure}}\par
\textbf{DOK rationale:} Part (c) coordinates sample structure, inferential target, hypotheses, shared expected counts, and conditions to adjudicate a procedure claim and explain why identical arithmetic does not imply identical population meaning.\par

\frameworkphaseheader{Do Now (first take) $\rightarrow$ Explore (video + follow-along) $\rightarrow$ Exit (finish + turn in)}{1 $\rightarrow$ 3}{45}{%
\begin{itemize}[leftmargin=1.5em,itemsep=2pt,topsep=2pt]
  \item Display the board and ask for one sentence using the study description.
  \item Begin the 8.5 video follow-along at minute 5.
  \item Return to the board at minute 33. Students finish the shortened parts and justify their judgment.
  \item Collect at the door. Score part (c) only using the three required elements.
\end{itemize}
}{%
\begin{itemize}[leftmargin=1.5em,itemsep=2pt,topsep=2pt]
  \item Write a one-sentence first take before the video.
  \item Complete the video follow-along worksheet.
  \item Finish (a)--(c), revise the first take in the exit line, and turn in the sheet.
\end{itemize}
}{%
\begin{itemize}[leftmargin=1.5em,itemsep=2pt,topsep=2pt]
  \item Which specific number or study detail supports your choice?
  \item What claim would go beyond the evidence, and why?
\end{itemize}
}{Circulate and ask for evidence. Accept a concise response; do not require omitted calculations or a second full inference write-up.}

\begin{callout}{\IconWarn\ WATCH FOR}{calloutred}
\begin{itemize}[leftmargin=1.5em,itemsep=2pt,topsep=2pt]
  \item Choosing a speaker without a reason.
  \item Copying numbers without explaining what they support.
\end{itemize}
\end{callout}

\sectionbanner{SCORING PART (c) --- E / P / I}

\textbf{Expected elements}
\begin{itemize}[leftmargin=1.5em,itemsep=2pt,topsep=2pt]
  \item \textbf{design-a} (required) --- Names homogeneity, cites separate route samples, and states equal ticket-method distributions
  \item \textbf{design-b} (required) --- Names independence, cites one two-variable sample, and states no association in the system population
  \item \textbf{judges-rule} (required) --- Explains why matching counts do not identify the sampling design or test
\end{itemize}

{\small\renewcommand{\arraystretch}{1.25}
\noindent\begin{tabular}{|p{0.5in}|p{5.9in}|}
\hline
\textbf{E} & Justifies both procedures and contextual null hypotheses, and explains why the table-only rule fails. \\ \hline
\textbf{P} & Shows substantive valid reasoning for at least one required element, but does not meet all E requirements. Credit correct reasoning even when another claim is wrong. \\ \hline
\textbf{I} & Shows no substantive valid reasoning for any required element; a name, unsupported choice, or copied number alone does not count. \\ \hline
\end{tabular}}

\clearpage
\focusbanner{TODAY'S PROBLEM --- annotated}

Suppose a transit authority obtains the shown counts under either design.\par\smallskip \textbf{Design A:} Take independent SRSs of 100 from each of three route populations of 5,000 and record ticket method.\par \textbf{Design B:} Take one SRS of 300 from 15,000 system riders and record both route and ticket method.\par\smallskip \textbf{Ari:} ``Route rows imply homogeneity; the counts are enough to choose.''\par \textbf{Bea:} ``Read how riders were sampled; identical counts can support different population claims.''\par
\begin{center}
\textbf{Ticket method}\par\smallskip
\begin{tabular}{|l|c|c|c|c|}
\hline
\textbf{Route} & \textbf{Mobile} & \textbf{Card} & \textbf{Cash} & \textbf{Total} \\ \hline
\textbf{North} & 60 & 25 & 15 & 100 \\ \hline
\textbf{East} & 48 & 32 & 20 & 100 \\ \hline
\textbf{West} & 42 & 33 & 25 & 100 \\ \hline
\textbf{Total} & 150 & 90 & 60 & 300 \\ \hline
\end{tabular}
\end{center}
\begin{firsttakebox}{FIRST TAKE (5 min)}
Do the same counts mean that riders were selected in the same way? Explain in one sentence.\par
\answer{Ungraded. Everyday language is enough before the video; formal vocabulary belongs in the finish.}
\end{firsttakebox}

\textit{Rules callout ``LET DATA COLLECTION NAME THE TEST (keep on the page)'' is printed on the student sheet.}\par\medskip

\dokbadge{1}~\textbf{(a)} How many random samples are taken in Design A and in Design B?\par
\answer{Design A takes three independent random samples; Design B takes one random sample.}

\dokbadge{2}~\textbf{(b)} Calculate the three expected counts for the North row. Explain briefly why the other two rows have the same expected counts.\par
\answer{North expected counts are $(100)(150)/300=50$, $(100)(90)/300=30$, and $(100)(60)/300=20$. The other rows share the same row total of 100 and column totals, so their expected counts are also $(50,30,20)$.}

\dokbadge{3}~\textbf{(c)} In 4 sentences, judge Ari's rule. Name and justify the test for each design, state each null hypothesis in context, and explain why the same table can lead to different test names.\par
\answer{Bea is justified: Design A uses homogeneity because it samples each route population separately. Its null says ticket-method distributions are the same across the three route populations. Design B uses independence because one system sample records two variables; its null says route and ticket method are not associated among system riders. Ari's row-label rule fails because test choice depends on data collection, even though the observed table and expected counts are identical.}

\begin{summaryexitbox}{EXIT}
One thing the video changed about my first take: \hrulefill
\end{summaryexitbox}

\end{document}
